% Báo cáo kỹ thuật WikiPulse -- bản cập nhật liên tục của cả nhóm.
% Định dạng theo mẫu báo cáo của Trường CNTT&TT, ĐHBK Hà Nội (lớp report 12pt, lề 2,5 cm, giãn dòng 1,5,
% phông Latin Modern, tài liệu tham khảo BibTeX kiểu plain). Quy ước cập nhật: xem chương đầu và README.md.
% Biên dịch: latexmk -xelatex main.tex  (hoặc xelatex main -> bibtex main -> xelatex main -> xelatex main)
\documentclass[12pt,a4paper]{report}

% ============================================================================
\usepackage{fontspec}               % XeLaTeX: phông Latin Modern (giống Computer Modern, đủ dấu tiếng Việt)
\usepackage[vietnamese]{babel}
\usepackage{amsmath,amssymb,amsfonts}
\usepackage{graphicx}
\usepackage{booktabs}
\usepackage{longtable}
\usepackage{array}
\usepackage{tabularx}
\usepackage{multirow}
\usepackage{algorithm}
\usepackage[noend]{algpseudocode}
\usepackage{listings}               % chỉ dùng bộ đếm và tên "Đoạn mã"; khối mã viết bằng fancyvrb (đúng chữ có dấu)
\usepackage{fancyvrb}
\usepackage[table]{xcolor}
\usepackage{hyperref}
\usepackage{geometry}
\usepackage{setspace}
\usepackage{caption}
\usepackage{subcaption}
\usepackage{float}
\usepackage{enumitem}
\usepackage{url}
\usepackage{tcolorbox}
\usepackage{tikz}
\usetikzlibrary{arrows.meta,positioning,shapes.geometric,shapes.misc,shapes.multipart,calc,fit,backgrounds}
\usepackage{pgf-umlsd}            % biểu đồ tuần tự UML
\usepackage{pgfplots}
\pgfplotsset{compat=1.18}

\geometry{margin=2.5cm}
\onehalfspacing
\graphicspath{{hinh/}}
\setcounter{tocdepth}{2}
\hypersetup{hidelinks}
\setlist{itemsep=1pt,topsep=3pt}
\emergencystretch=3em          % tránh tràn lề khi dòng có tên mã dài không ngắt được

\newcommand{\enquote}[1]{“#1”}
\newcommand{\code}[1]{\texttt{#1}}
\newcommand{\chualam}{\par\noindent\textit{(Dự kiến, chưa thực hiện.)}\par}
% Ghi chú nhỏ dưới một mục đã viết nhưng còn sửa ở các tuần sau
\newcommand{\dukiensua}[1]{\par\noindent{\footnotesize\itshape\color{black!60}Ghi chú: #1}\par}
% Thông tin còn thiếu, cần bổ sung
\newcommand{\canbosung}[1]{{\color{red}[#1]}}

\renewcommand{\lstlistingname}{Đoạn mã}

% Từ khóa giả mã giữ tiếng Anh chuẩn (Input, Output, for, if, return...); chỉ phần mô tả viết tiếng Việt
\algrenewcommand\algorithmicrequire{\textbf{Input:}}
\algrenewcommand\algorithmicensure{\textbf{Output:}}
\algnewcommand\algorithmicand{\textbf{and}}
\algnewcommand\AlgAnd{\algorithmicand}
\floatname{algorithm}{Thuật toán}

% Màu dùng trong sơ đồ
\definecolor{wpblue}{RGB}{18,92,173}
\definecolor{wpgreen}{RGB}{46,140,87}
\definecolor{wpamber}{RGB}{214,140,20}
\definecolor{wpred}{RGB}{190,50,50}
\definecolor{wpgray}{RGB}{110,116,128}
\definecolor{lightblue}{RGB}{228,239,252}
\definecolor{lightgreen}{RGB}{224,244,232}
\definecolor{lightamber}{RGB}{255,241,214}
\definecolor{lightred}{RGB}{252,226,226}
\definecolor{lightgray}{RGB}{242,243,245}

% Hộp ghi chú: khung mảnh, nền xám nhạt, hợp văn phong báo cáo
\newtcolorbox{ghichu}[1]{colback=lightgray,colframe=black!55,boxrule=0.4pt,arc=0pt,
  fonttitle=\bfseries,coltitle=black,colbacktitle=black!10,title=#1}

\tikzset{
  box/.style={draw=wpblue,fill=lightblue,rounded corners=3pt,align=center,font=\small\sffamily,
              minimum height=9mm,inner sep=4pt,text width=#1},
  box/.default=28mm,
  gbox/.style={box=#1,draw=wpgreen,fill=lightgreen},
  gbox/.default=28mm,
  abox/.style={box=#1,draw=wpamber,fill=lightamber},
  abox/.default=28mm,
  rbox/.style={box=#1,draw=wpred,fill=lightred},
  rbox/.default=28mm,
  db/.style={cylinder,shape border rotate=90,aspect=0.18,draw=wpblue,fill=lightblue,
             align=center,font=\small\sffamily,minimum height=13mm,text width=24mm},
  dec/.style={diamond,aspect=2.2,draw=wpamber,fill=lightamber,align=center,font=\small\sffamily,inner sep=1pt},
  arr/.style={-{Stealth[length=2.6mm]},thick,color=wpgray},
  lab/.style={font=\scriptsize\sffamily,color=wpgray,fill=white,inner sep=1pt},
  uc/.style={ellipse,draw=wpblue,fill=lightblue,align=center,font=\footnotesize\sffamily,
             minimum width=40mm,minimum height=9mm,inner sep=1pt},
  actor/.pic={\draw[thick] (0,0.42) circle (0.16);
              \draw[thick] (0,0.26) -- (0,-0.12) (-0.28,0.14) -- (0.28,0.14) (0,-0.12) -- (-0.2,-0.46) (0,-0.12) -- (0.2,-0.46);},
  % --- UML: biểu đồ hoạt động ---
  start/.style={circle,fill=black,minimum size=3.2mm,inner sep=0},
  stop/.style={circle,draw=black,thick,fill=white,minimum size=4.4mm,inner sep=0,
               path picture={\fill[black] (path picture bounding box.center) circle (1.25mm);}},
  act/.style={rectangle,rounded corners=2.5mm,draw=black,fill=lightblue,align=center,font=\footnotesize\sffamily,
              minimum height=8mm,inner sep=3pt,text width=#1},
  act/.default=40mm,
  dnode/.style={diamond,draw=black,fill=white,minimum size=6.5mm,inner sep=0},
  flow/.style={-{Stealth[length=2.2mm]},thick},
  guard/.style={font=\scriptsize\sffamily,fill=white,inner sep=1pt},
  % --- UML: biểu đồ trạng thái ---
  state/.style={rectangle,rounded corners=3mm,draw=black,fill=lightblue,align=center,font=\footnotesize\sffamily,
                minimum height=9mm,minimum width=24mm,inner sep=3pt},
  % --- UML: biểu đồ gói, triển khai ---
  pkg/.style={rectangle,draw=black,fill=lightblue,align=center,font=\footnotesize\sffamily,minimum width=24mm,
              minimum height=9mm,inner sep=2pt,
              append after command={\pgfextra{\draw[fill=lightblue] (\tikzlastnode.north west) rectangle ++(8mm,2.4mm);}}},
  dep/.style={-{Straight Barb[length=2mm]},dashed,thick},
  udnode/.style={rectangle,draw=black,fill=lightblue,align=center,font=\footnotesize\sffamily,inner sep=4pt,
                 append after command={\pgfextra{\draw (\tikzlastnode.north west) -- ++(3mm,3mm)
                   -- ($(\tikzlastnode.north east)+(3mm,3mm)$) -- ($(\tikzlastnode.south east)+(3mm,3mm)$)
                   -- (\tikzlastnode.south east) ($(\tikzlastnode.north east)+(3mm,3mm)$) -- (\tikzlastnode.north east);}}},
  % --- sơ đồ CSDL (sinh bởi tools/erd.py) ---
  erd/.style={rectangle split,rectangle split parts=2,draw=black,anchor=north west,font=\scriptsize\sffamily,
              rectangle split part fill={lightblue,white},inner sep=2pt},
  erdstub/.style={draw=black,dashed,anchor=north west,font=\scriptsize\sffamily,inner sep=4pt,fill=lightgray},
  erdline/.style={thick},
  erdmult/.style={font=\scriptsize\sffamily,fill=white,inner sep=1pt},
  erdfk/.style={font=\tiny\sffamily,color=wpgray,sloped,above,inner sep=1pt},
}

% ============================================================================
% DOCUMENT BEGIN
% ============================================================================

% ============================================================================
% DOCUMENT BEGIN
% ============================================================================
\begin{document}

% ============================================================================
% TRANG BÌA (theo mẫu báo cáo của Trường CNTT&TT)
% ============================================================================
\begin{titlepage}
\begin{center}

\textsc{\large \textbf{ĐẠI HỌC BÁCH KHOA HÀ NỘI}}\\[0.5cm]

{\Large TRƯỜNG CÔNG NGHỆ THÔNG TIN VÀ TRUYỀN THÔNG \\[0.6cm]}

\includegraphics[scale=0.11]{logo-soict-hust-1.png}\\[0.6cm]
\textsc{\Large \textbf{\canbosung{Mã học phần -- Tên học phần}}}\\[0.4cm]

\rule{\linewidth}{0.7mm} \\[0.4cm]
{\huge \bfseries WikiPulse\\[0.3cm]}
{\Large \bfseries Phân tích xu hướng Wikipedia và ưu tiên xem xét\\[0.1cm] chỉnh sửa đáng ngờ theo thời gian thực\\[0.3cm]}
{\large (Báo cáo kỹ thuật -- bản cập nhật liên tục)\\[0.3cm]}
\rule{\linewidth}{0.7mm} \\[0.6cm]

{\large \textbf{Giảng viên hướng dẫn:}}\\[0.3cm]
\begin{tabular}{l}
\large \canbosung{Học hàm, học vị, họ tên} \\[0.4cm]
\end{tabular}

{\large \textbf{Nhóm sinh viên thực hiện:}}\\[0.3cm]
\begin{tabular}{ll}
\large \canbosung{Họ tên} (Trọng) & \large \canbosung{MSSV} \\
\large \canbosung{Họ tên} (Thành) & \large \canbosung{MSSV} \\
\large \canbosung{Họ tên} (Linh) & \large \canbosung{MSSV} \\
\large Hoàng Văn Nhân & \large 20235542 \\
\large \canbosung{Họ tên} (Mạnh) & \large \canbosung{MSSV} \\
\end{tabular}

\vfill

{\large \textbf{Hà Nội - Việt Nam}\\[0.2cm] \color{blue!90!black}2026}

\end{center}
\end{titlepage}
\setstretch{1.5}

\pagenumbering{roman}

\chapter*{Quy ước cập nhật báo cáo}
\addcontentsline{toc}{chapter}{Quy ước cập nhật báo cáo}
Báo cáo này được cả nhóm viết dần trong quá trình phát triển, để thành viên mới đọc hiểu được project và để báo cáo
cuối kỳ không phải soạn lại từ đầu. Quy ước:
\begin{itemize}
  \item Phần chưa làm chỉ có đề mục, kèm dòng \textit{(Dự kiến, chưa thực hiện.)} -- macro \code{\textbackslash chualam}.
  \item Phần đã viết nhưng còn sửa theo các tuần sau có ghi chú nhỏ bên dưới -- macro
        \code{\textbackslash dukiensua\{nội dung, tuần\}}.
  \item Thông tin còn thiếu đánh dấu màu đỏ -- macro \code{\textbackslash canbosung}. Trên trang bìa, mỗi thành viên tự
        điền họ tên và MSSV của mình; trưởng nhóm điền học phần và giảng viên.
  \item Mỗi việc trong \code{docs/checkpoints.md} khi chuyển sang DONE thì owner cập nhật mục tương ứng ở đây trong cùng
        PR, kèm số liệu và bằng chứng tái hiện được. Không ghi kết quả chưa đo.
  \item Nguồn sự thật cho giao diện giữa các module vẫn là \code{docs/contracts.md}; báo cáo giải thích và dẫn lại, không
        định nghĩa contract mới.
\end{itemize}

\begin{table}[H]
\centering\small
\caption{Lịch sử cập nhật báo cáo}\label{tab:lichsu}
\begin{tabularx}{\linewidth}{@{}llX@{}}
\toprule
\textbf{Ngày} & \textbf{Người} & \textbf{Nội dung} \\ \midrule
15/10/2026 & Nhân & Khởi tạo khung báo cáo; viết phần CP0 (kiến trúc, luồng dữ liệu, serving, kiểm thử) \\
15/10/2026 & Nhân & Thêm use case (biểu đồ, đặc tả), biểu đồ hoạt động; viết lại Chương 2 kèm hình minh họa; chuyển
tên file, lệnh sang Phụ lục A \\
\bottomrule
\end{tabularx}
\end{table}

\tableofcontents
\newpage

\begin{abstract}
Wikipedia nhận hàng trăm chỉnh sửa mỗi phút. Wikimedia công bố toàn bộ thay đổi này dưới dạng luồng sự kiện công khai
(EventStreams \code{recentchange})~\cite{eventstreams}. Luồng này đủ để biết trang nào đang được quan tâm và chỉnh sửa nào
đáng được người tuần tra xem trước, nhưng dữ liệu đến liên tục, có trùng lặp, có sự kiện đến muộn và có bản ghi hỏng.

Báo cáo trình bày \textbf{WikiPulse}, một nền tảng xử lý luồng đầy đủ: Kafka nhận sự kiện, Spark Structured Streaming
kiểm tra, chuẩn hóa, loại trùng và tổng hợp theo cửa sổ thời gian sự kiện, dữ liệu được lưu ba tầng Bronze/Silver/Gold
dạng Parquet, kết quả phục vụ qua Cassandra, FastAPI và dashboard Streamlit. Hệ thống chạy được ở hai chế độ: LIVE nối
với Wikimedia, và REPLAY phát lại dữ liệu mẫu cố định để kiểm thử tái hiện được.

Đến ngày 15/10/2026, nhóm đã hoàn thành mốc CP0: luồng REPLAY chạy thật từ Kafka tới dashboard cho nghiệp vụ đếm hoạt
động theo wiki trong cửa sổ 1 phút, kết quả khớp bộ kết quả kỳ vọng độc lập và được kiểm chứng trên hai máy (macOS ARM64
và Windows x86\_64). Trang xu hướng, ưu tiên chỉnh sửa đáng ngờ, chế độ LIVE và HDFS là việc của tuần 2--4, được liệt
kê dưới dạng đề mục.
\dukiensua{Tóm tắt viết lại khi chốt CP4.}
\end{abstract}

\chapter*{Danh mục thuật ngữ và viết tắt}
\addcontentsline{toc}{chapter}{Danh mục thuật ngữ và viết tắt}
\begin{tabularx}{\linewidth}{@{}lX@{}}
\toprule
\textbf{Thuật ngữ} & \textbf{Nghĩa} \\ \midrule
Event time & Thời điểm sự kiện xảy ra trên Wikipedia (trường \code{meta.dt}), khác thời điểm hệ thống nhận \\
Cửa sổ (window) & Khoảng thời gian event time cố định để tổng hợp, ví dụ 1 phút \\
Watermark & Ngưỡng trễ tối đa hệ thống chờ sự kiện đến muộn (10 phút); qua ngưỡng thì cửa sổ được chốt \\
Bronze / Silver / Gold & Ba tầng dữ liệu: thô, đã chuẩn hóa, đã tổng hợp và chốt \\
Provisional / finalized & Số liệu tạm tính (có thể còn đổi) / số liệu đã chốt sau watermark \\
REPLAY / LIVE & Chế độ phát lại dữ liệu mẫu / chế độ nhận luồng thật từ Wikimedia \\
Upsert & Ghi đè theo khóa: có thì thay, chưa có thì thêm \\
DLQ & Dead-letter queue, topic chứa bản ghi không xử lý được \\
Canary & Sự kiện kiểm tra do Wikimedia phát vào luồng, không phải chỉnh sửa thật \\
Checkpoint & Trạng thái Spark lưu lại để chạy tiếp sau khi dừng; đồng thời là mốc nghiệm thu CP0--CP4 của nhóm \\
\bottomrule
\end{tabularx}

\newpage
\pagenumbering{arabic}

% ============================================================================
% CHƯƠNG 1: GIỚI THIỆU
% ============================================================================
\chapter{Giới thiệu}\label{chap:gioithieu}

\section{Đặt vấn đề}
Luồng \code{recentchange} của Wikimedia phát mọi thay đổi trên các dự án Wikimedia: sửa trang, tạo trang mới, thay đổi
nhật ký. Mỗi sự kiện có wiki, tiêu đề trang, người sửa, kích thước trước và sau, cờ bot, cờ ẩn danh. Từ luồng này có thể
trả lời hai câu hỏi có ích:
\begin{enumerate}
  \item \textbf{Đang có gì nóng?} Trang nào đang được sửa nhiều bất thường trong vài phút gần nhất.
  \item \textbf{Nên xem chỉnh sửa nào trước?} Chỉnh sửa nào có dấu hiệu cần người tuần tra kiểm tra, ví dụ xóa nhiều
        nội dung, hoàn tác qua lại, hay sửa ẩn danh.
\end{enumerate}

Khó khăn nằm ở tính chất của dữ liệu luồng: sự kiện đến không theo thứ tự, có bản trùng khi nguồn kết nối lại, có bản
ghi hỏng, và kết quả phải cập nhật liên tục trong khi vẫn phải đúng khi hệ thống khởi động lại.

\section{Mục tiêu}
\begin{itemize}
  \item Xây dựng luồng end-to-end Kafka $\rightarrow$ Spark $\rightarrow$ Parquet/Cassandra $\rightarrow$ API
        $\rightarrow$ dashboard, chạy được ở hai chế độ LIVE và REPLAY.
  \item Tổng hợp theo event time có watermark, loại trùng, cách ly bản ghi lỗi, đúng khi retry và khởi động lại.
  \item Phục vụ ba nghiệp vụ: hoạt động theo wiki, trang xu hướng, danh sách chỉnh sửa cần ưu tiên xem xét.
  \item Bảo vệ quyền riêng tư: tên người dùng và IP không đi ra khỏi tầng Bronze.
\end{itemize}

\section{Phạm vi}
\begin{itemize}
  \item \textbf{Trong phạm vi:} MVP một node chạy bằng Docker Compose; mặc định hai wiki \code{enwiki}, \code{viwiki}, mở
        rộng bằng cấu hình; lưu Parquet cục bộ (tuần 1) rồi HDFS (tuần 2--3).
  \item \textbf{Ngoài phạm vi:} mô hình học máy, độ sẵn sàng cao, triển khai cloud, xác thực và phân quyền. Kubernetes và
        monitoring chỉ ở mức khung.
  \item \textbf{Nguyên tắc:} điểm rủi ro chỉ dùng để \emph{ưu tiên xem xét}, không kết luận một chỉnh sửa là phá hoại.
        Không đưa dữ liệu giả vào đường demo.
\end{itemize}

\section{Tổ chức nhóm}
Nhóm 5 người, 4 tuần, khoảng 10--15 giờ mỗi người mỗi tuần. Tiến độ quản lý theo mốc nghiệm thu
(Bảng~\ref{tab:moc}) và bảng việc chi tiết trong \code{docs/checkpoints.md}; một việc chỉ chuyển DONE khi có bằng chứng
kiểm thử.

\begin{table}[H]
\centering\small
\caption{Phân công theo module}\label{tab:phancong}
\begin{tabularx}{\linewidth}{@{}lX@{}}
\toprule
\textbf{Thành viên} & \textbf{Phụ trách} \\ \midrule
Trọng & Core, cấu hình, Docker Compose, Makefile, tích hợp và nghiệm thu chung \\
Thành & Ingestion: Kafka topics, replay, collector LIVE \\
Linh & Xử lý Spark: validation, chuẩn hóa, loại trùng, cửa sổ, đặc trưng xu hướng và rủi ro \\
Nhân & Lưu trữ phục vụ Cassandra, sink từ Spark, API FastAPI \\
Mạnh & Dashboard Streamlit, kiểm thử chấp nhận (acceptance) \\
\bottomrule
\end{tabularx}
\end{table}

\begin{table}[H]
\centering\small
\caption{Các mốc nghiệm thu}\label{tab:moc}
\begin{tabularx}{\linewidth}{@{}lXl@{}}
\toprule
\textbf{Mốc} & \textbf{Điều kiện DONE} & \textbf{Trạng thái} \\ \midrule
CP0 & REPLAY thật qua Kafka/Spark/3 tầng Parquet/Cassandra/API/Overview & DONE (11/10/2026) \\
CP1 & Người khác chạy lại được CP0; mỗi người nhận module, contract và test đầu tiên & Đang làm \\
CP2 & LIVE/REPLAY, cửa sổ, rủi ro, xu hướng đúng fixture; UI thật; HDFS khởi động & Chưa \\
CP3 & E2E trên HDFS; late data, khởi động lại, retry, Gold finalized được kiểm chứng & Chưa \\
CP4 & Clone sạch trên máy thứ hai; replay pass; LIVE $\geq$ 15 phút; 4 trang; đủ tài liệu & Chưa \\
\bottomrule
\end{tabularx}
\end{table}
\dukiensua{Cập nhật cột trạng thái mỗi khi đạt mốc.}

\section{Cấu trúc báo cáo}
\begin{itemize}
  \item Chương~\ref{chap:nen} trình bày kiến thức nền về xử lý luồng và các công nghệ sử dụng.
  \item Chương~\ref{chap:lienquan} điểm lại các công cụ và nghiên cứu liên quan.
  \item Chương~\ref{chap:thietke} mô tả phân tích, thiết kế và cài đặt.
  \item Chương~\ref{chap:danhgia} trình bày kiểm thử và kết quả.
  \item Chương~\ref{chap:ketluan} kết luận và liệt kê phần còn lại.
\end{itemize}

% ============================================================================
% CHƯƠNG 2: CƠ SỞ LÝ THUYẾT VÀ CÔNG NGHỆ
% ============================================================================
\chapter{Cơ sở lý thuyết và công nghệ}\label{chap:nen}

Chương này giải thích các khái niệm mà thiết kế ở Chương~\ref{chap:thietke} dựa vào: dữ liệu luồng khác dữ liệu tĩnh ở
đâu, Kafka giữ sự kiện thế nào, Spark tổng hợp theo thời gian ra sao khi sự kiện đến lộn xộn, vì sao dữ liệu được lưu
thành ba tầng, và vì sao Cassandra được chọn để phục vụ kết quả.

\section{Bài toán xử lý dữ liệu luồng}\label{sec:luong}
Dữ liệu tĩnh (một file, một bảng) có điểm bắt đầu và kết thúc: chạy truy vấn một lần là có kết quả. Dữ liệu luồng thì
\textbf{không bao giờ kết thúc}: sự kiện đến liên tục, kết quả phải được tính dần và cập nhật mãi. Ngoài ra, mỗi sự kiện
có hai mốc thời gian khác nhau:
\begin{itemize}
  \item \textbf{Event time}: lúc sự việc thật sự xảy ra, ví dụ lúc một người bấm lưu chỉnh sửa trên Wikipedia.
  \item \textbf{Processing time}: lúc hệ thống nhận được sự kiện đó.
\end{itemize}
Hai mốc này lệch nhau vì mạng chậm, nguồn mất kết nối rồi gửi bù, hay hệ thống tạm dừng. Hình~\ref{fig:eventtime} cho thấy
ba hiện tượng mà mọi hệ thống luồng phải xử lý.

\begin{figure}[H]
\centering
\begin{tikzpicture}[font=\footnotesize\sffamily]
  \draw[arr] (0,2.2) -- (12.4,2.2) node[right,align=left]{event time\\(lúc xảy ra)};
  \draw[arr] (0,0) -- (12.4,0) node[right,align=left]{processing time\\(lúc nhận)};
  \foreach \x/\l in {0.5/00:00,3.5/00:01,6.5/00:02,9.5/00:03} {
    \draw (\x,2.1) -- (\x,2.3); \node[wpgray] at (\x,3.05) {\l}; \draw[wpgray!50,densely dotted] (\x,2.3) -- (\x,2.85);
    \draw (\x,-0.1) -- (\x,0.1);
  }
  % đúng thứ tự
  \foreach \a/\b/\n in {1.0/1.6/e1, 2.4/3.0/e2} {
    \fill[wpblue] (\a,2.2) circle (0.09) node[above=1mm,text=wpblue]{\n};
    \draw[wpblue,thick,-{Stealth[length=2mm]}] (\a,2.1) -- (\b,0.1);
    \fill[wpblue] (\b,0) circle (0.09);
  }
  % đến muộn, lộn thứ tự
  \fill[wpamber] (3.0,2.2) circle (0.09) node[above=1mm,text=wpamber]{e3};
  \draw[wpamber,thick,-{Stealth[length=2mm]}] (3.0,2.1) -- (7.6,0.1);
  \fill[wpamber] (7.6,0) circle (0.09);
  \fill[wpblue] (4.6,2.2) circle (0.09) node[above=1mm,text=wpblue]{e4};
  \draw[wpblue,thick,-{Stealth[length=2mm]}] (4.6,2.1) -- (5.2,0.1);
  \fill[wpblue] (5.2,0) circle (0.09);
  % trùng
  \fill[wpred] (7.0,2.2) circle (0.09) node[above=1mm,text=wpred]{e5};
  \draw[wpred,thick,-{Stealth[length=2mm]}] (7.0,2.1) -- (8.6,0.1);
  \draw[wpred,thick,dashed,-{Stealth[length=2mm]}] (7.0,2.1) -- (10.4,0.1);
  \fill[wpred] (8.6,0) circle (0.09); \fill[wpred] (10.4,0) circle (0.09);
  % hỏng
  \node[wpgray,font=\footnotesize\bfseries] at (11.5,0) {\{\,?\,\}};
  \node[below,lab,text=wpamber] at (7.6,-0.25) {e3 đến sau e4};
  \node[below,lab,text=wpred,align=center] at (9.5,-0.25) {e5 đến hai lần};
  \node[below,lab,align=center] at (11.6,-0.25) {bản ghi hỏng};
\end{tikzpicture}
\caption{Event time và processing time. Mỗi mũi tên nối lúc sự kiện xảy ra với lúc hệ thống nhận được. Sự kiện có thể đến
\textcolor{wpamber}{muộn và lộn thứ tự} (e3), \textcolor{wpred}{đến hai lần} khi nguồn gửi lại (e5), hoặc đến ở dạng
hỏng không đọc được.}\label{fig:eventtime}
\end{figure}

Nếu đếm theo processing time, kết quả phụ thuộc vào việc mạng nhanh hay chậm, và chạy lại cùng dữ liệu sẽ ra số khác.
Vì vậy WikiPulse đếm theo \textbf{event time}: chạy lại bao nhiêu lần cũng ra cùng một kết quả. Cái giá phải trả là hệ
thống phải chờ sự kiện muộn, phải tự loại bản trùng và tách riêng bản ghi hỏng. Các mục sau giải thích cách làm từng việc.

\section{Luồng sự kiện thay đổi của Wikipedia}
Wikimedia công bố mọi thay đổi trên các dự án của mình qua dịch vụ EventStreams~\cite{eventstreams}. Client mở một kết nối
HTTP dài và nhận từng sự kiện dạng JSON ngay khi nó xảy ra (giao thức Server-Sent Events). Bảng~\ref{tab:recentchange} liệt
kê các thông tin của một sự kiện mà hệ thống dùng tới.

\begin{table}[H]
\centering\small
\caption{Thông tin chính trong một sự kiện thay đổi và cách hệ thống dùng}\label{tab:recentchange}
\begin{tabularx}{\linewidth}{@{}lXX@{}}
\toprule
\textbf{Thông tin} & \textbf{Ý nghĩa} & \textbf{Dùng để} \\ \midrule
Định danh sự kiện & Duy nhất cho mỗi sự kiện & Loại bản trùng \\
Thời điểm & Lúc thay đổi xảy ra (event time) & Xếp vào cửa sổ thời gian \\
Wiki, trang & Ví dụ \enquote{enwiki}, tên bài viết & Nhóm kết quả theo wiki, theo trang \\
Loại thay đổi & Sửa trang, tạo trang mới, ghi nhật ký & Đếm từng loại \\
Kích thước trước, sau & Số byte của trang & Phát hiện xóa nhiều nội dung \\
Người sửa, cờ ẩn danh, cờ bot & Tài khoản hoặc địa chỉ IP & Đếm sửa ẩn danh, sửa bởi bot; được băm trước khi lưu \\
\bottomrule
\end{tabularx}
\end{table}

Kết nối dài có thể đứt. Khi nối lại, client có thể nhận lại một phần sự kiện đã nhận, nên phía sau phải chịu được bản
trùng. Ngoài sự kiện thật, Wikimedia còn phát định kỳ \textbf{sự kiện canary} để tự kiểm tra hệ thống của họ; các sự
kiện này phải bị loại khỏi số liệu.
\dukiensua{Bổ sung cách nối lại không mất sự kiện khi làm collector LIVE (W2-THANH-01).}

\section{Apache Kafka: bộ đệm sự kiện bền vững}
Kafka~\cite{kreps2011} lưu sự kiện thành một \textbf{nhật ký chỉ ghi thêm} (append-only log). Các khái niệm chính
(Hình~\ref{fig:kafka}):
\begin{itemize}
  \item \textbf{Topic}: một kênh sự kiện, ví dụ kênh chứa các thay đổi thô của Wikipedia.
  \item \textbf{Partition}: topic được chia thành nhiều nhật ký song song để chia tải. Mỗi sự kiện trong một partition có
        số thứ tự tăng dần gọi là \textbf{offset}.
  \item \textbf{Khóa}: producer gắn khóa cho mỗi sự kiện; sự kiện cùng khóa luôn vào cùng partition, nên giữ đúng thứ tự
        với nhau. WikiPulse dùng khóa là wiki và trang, để các chỉnh sửa của cùng một trang không bị đảo thứ tự.
  \item \textbf{Consumer} tự nhớ mình đã đọc đến offset nào, nên có thể dừng rồi đọc tiếp, hoặc đọc lại từ đầu.
\end{itemize}

\begin{figure}[H]
\centering
\begin{tikzpicture}[font=\footnotesize\sffamily]
  \node[box=24mm] (p1) at (-0.6,1.5) {Replay\\(dữ liệu mẫu)};
  \node[box=24mm,dashed] (p2) at (-0.6,-1.5) {Collector LIVE};
  \foreach \r/\y/\n in {0/1.5/5, 1/0/4, 2/-1.5/6} {
    \node[font=\scriptsize\sffamily,anchor=east] at (3.2,\y) {partition \r};
    \foreach \i in {0,...,\n} {
      \node[draw=wpblue,fill=lightblue,minimum width=6mm,minimum height=6mm,inner sep=0,font=\scriptsize\sffamily]
        at (3.6+\i*0.62,\y) {\i};
    }
  }
  \node[draw=black,dashed,fit={(1.9,1.9) (8.0,-1.9)},inner sep=2mm,
        label={[font=\footnotesize\sffamily\bfseries]above:topic: các thay đổi thô}] (topic) {};
  \node[box=24mm] (c) at (11.4,0) {Spark\\(consumer)};
  \draw[arr] (p1.east) -- (1.95,1.5);
  \draw[arr,dashed] (p2.east) -- (1.95,-1.5);
  \draw[arr] (topic.east) -- node[above,lab]{đọc theo offset} (c.west);
  \draw[wpred,thick,-{Stealth[length=2mm]}] (5.46,-0.75) -- (5.46,-0.35);
  \node[font=\scriptsize\sffamily,text=wpred,anchor=west] at (5.6,-0.62) {đã đọc tới đây};
  \node[font=\scriptsize\sffamily,text=wpgray] at (5.0,-2.55) {sự kiện mới luôn ghi vào cuối $\rightarrow$};
\end{tikzpicture}
\caption{Topic Kafka gồm ba partition. Mỗi ô là một sự kiện, số trong ô là offset. Producer ghi vào cuối nhật ký;
consumer (Spark) đọc theo offset và nhớ vị trí đã đọc. Nét đứt là phần chưa làm.}\label{fig:kafka}
\end{figure}

Trong WikiPulse, Kafka đứng giữa nguồn sự kiện và Spark. Spark dừng để nâng cấp hay bị lỗi thì sự kiện vẫn nằm an toàn
trong Kafka; khi Spark chạy lại, nó đọc tiếp từ offset cũ. Nhờ đọc lại được, chế độ REPLAY chỉ cần ghi dữ liệu mẫu vào
cùng topic là toàn bộ phần sau chạy y như với dữ liệu thật.

\section{Spark Structured Streaming}
\subsection{Mô hình micro-batch}
Spark Structured Streaming~\cite{armbrust2018} cho phép viết truy vấn trên luồng giống như truy vấn trên một bảng bình
thường. Luồng được coi là một \textbf{bảng không giới hạn}, mỗi sự kiện mới là một dòng thêm vào. Cứ sau một khoảng cố
định (10 giây trong WikiPulse), Spark lấy phần dữ liệu mới thành một \textbf{micro-batch}, tính tiếp kết quả từ trạng thái
lần trước, rồi ghi kết quả ra đích (Hình~\ref{fig:microbatch}).

\begin{figure}[H]
\centering
\begin{tikzpicture}[font=\footnotesize\sffamily]
  \foreach \i/\x in {1/0,2/2.3,3/4.6} {
    \node[draw=wpblue,fill=lightblue,minimum width=19mm,minimum height=8mm] (b\i) at (\x,0) {batch \i};
  }
  \node[font=\footnotesize] at (6.4,0) {\ldots};
  \draw[arr] (-1.3,-0.9) -- (7.0,-0.9) node[right]{thời gian};
  \foreach \x/\t in {-0.95/0s,1.35/10s,3.65/20s,5.95/30s} \draw (\x,-0.8) -- (\x,-1.0) node[below]{\t};
  \node[box=30mm] (q) at (2.3,-2.8) {Truy vấn\\(lọc, chuẩn hóa, đếm)};
  \node[db] (st) at (-2.4,-2.8) {Checkpoint\\offset + trạng thái};
  \node[db] (out) at (6.2,-2.8) {Đích ghi\\(Parquet, Cassandra)};
  \draw[arr] (b3.south) -- (q.north);
  \draw[arr,<->] (st) -- node[above,lab]{đọc, lưu} (q);
  \draw[arr] (q) -- (out);
\end{tikzpicture}
\caption{Mô hình micro-batch. Mỗi 10 giây, phần dữ liệu mới thành một batch; truy vấn kết hợp batch với trạng thái lưu
trong checkpoint để ra kết quả mới.}\label{fig:microbatch}
\end{figure}

\textbf{Checkpoint} lưu hai thứ: đã đọc Kafka tới offset nào, và trạng thái đang tính dở (ví dụ số đếm của các cửa sổ chưa
chốt). Khi Spark khởi động lại, nó đọc checkpoint và chạy tiếp như chưa từng dừng. Checkpoint gắn chặt với cấu trúc
truy vấn: đổi cách tính có trạng thái thì phải dùng checkpoint mới.

\subsection{Cửa sổ thời gian và watermark}
Để đếm \enquote{mỗi phút có bao nhiêu chỉnh sửa}, sự kiện được xếp vào các \textbf{cửa sổ} 1 phút theo event time. Câu hỏi
khó là: \emph{khi nào thì một cửa sổ được coi là đủ?} Sự kiện của phút 00:00 có thể vẫn đang trên đường tới. Spark trả
lời bằng \textbf{watermark}~\cite{akidau2015}:
\[
  \text{watermark} = \text{event time lớn nhất đã thấy} - \text{độ trễ cho phép (10 phút)}.
\]
Cửa sổ nào kết thúc trước watermark thì được \textbf{chốt}: kết quả không đổi nữa và trạng thái của nó được xóa khỏi bộ
nhớ. Sự kiện đến muộn hơn watermark bị bỏ khỏi phép đếm (nhưng vẫn được lưu ở tầng thô để tra lại).
Hình~\ref{fig:watermark} minh họa bằng chính dữ liệu mẫu của dự án.

\begin{figure}[H]
\centering
\begin{tikzpicture}[font=\footnotesize\sffamily,x=0.85cm]
  \fill[lightgreen] (0,-0.2) rectangle (2.17,1.6);
  \fill[lightamber] (2.17,-0.2) rectangle (14.5,1.6);
  \draw[arr] (0,0) -- (14.8,0) node[right]{event time};
  \foreach \m in {0,1,2,...,14} \draw (\m,-0.08) -- (\m,0.08);
  \foreach \m/\t in {0/00:00,2/00:02,4/00:04,6/00:06,8/00:08,10/00:10,12/00:12,14/00:14}
    \node[below] at (\m,-0.1) {\t};
  \draw[wpblue,thick] (0,0.35) -- (0,0.55) -- (1,0.55) -- (1,0.35);
  \node[above,text=wpblue] at (0.5,0.55) {cửa sổ 00:00};
  \foreach \x in {0.15,0.3,0.45,0.6,0.8} \fill[wpblue] (\x,0.2) circle (0.06);
  \fill[wpblue] (12.17,0.2) circle (0.07);
  \draw[wpgray,-{Stealth[length=2mm]}] (12.17,1.15) -- (12.17,0.3);
  \node[above,align=center,text=wpgray] at (12.17,1.1) {sự kiện mới nhất\\00:12:10};
  \draw[wpred,very thick,dashed] (2.17,-0.2) -- (2.17,1.6);
  \node[above,text=wpred,align=center] at (2.17,1.6) {watermark = 00:12:10 $-$ 10 phút = 00:02:10};
  \node[text=wpgreen!70!black,font=\footnotesize\bfseries] at (1.1,1.3) {đã chốt};
  \node[text=wpamber!80!black,font=\footnotesize\bfseries] at (7.5,1.3) {còn mở, số tạm tính};
  \fill[wpred] (0.7,-0.75) circle (0.07);
  \draw[wpred,thick] (0.55,-0.9) -- (0.85,-0.6) (0.55,-0.6) -- (0.85,-0.9);
  \node[right,text=wpred] at (1.0,-0.75) {sự kiện 00:00:40 đến lúc này: quá muộn, không được đếm};
\end{tikzpicture}
\caption{Watermark trên dữ liệu mẫu. Sự kiện lúc 00:12:10 đẩy watermark tới 00:02:10, nên cửa sổ 00:00 (5 chỉnh sửa
enwiki) được chốt. Mọi cửa sổ sau watermark vẫn mở và có thể còn thay đổi.}\label{fig:watermark}
\end{figure}

Độ trễ cho phép là một sự đánh đổi: dài thì đếm được nhiều sự kiện muộn hơn nhưng phải giữ nhiều trạng thái và chờ lâu
mới có số chốt; ngắn thì ngược lại. Dự án chọn 10 phút.

\subsection{Số tạm tính và số đã chốt}
Người xem dashboard muốn thấy số ngay, không muốn chờ 10 phút. Spark có hai \textbf{chế độ xuất} kết quả, và WikiPulse
dùng cả hai cho cùng một phép đếm (Bảng~\ref{tab:chedo}).

\begin{table}[H]
\centering\small
\caption{Hai chế độ xuất kết quả}\label{tab:chedo}
\begin{tabularx}{\linewidth}{@{}lXXl@{}}
\toprule
\textbf{Chế độ} & \textbf{Khi nào ghi} & \textbf{Tính chất} & \textbf{Đích trong WikiPulse} \\ \midrule
Update & Mỗi batch, mọi cửa sổ vừa thay đổi & Nhanh; số còn có thể tăng (tạm tính) & Cassandra, cho API \\
Append & Chỉ khi cửa sổ đã qua watermark & Chậm hơn ít nhất 10 phút; số không đổi nữa (đã chốt) & Tầng Gold, cho phân tích \\
\bottomrule
\end{tabularx}
\end{table}

Spark có một giới hạn quan trọng: một truy vấn ở chế độ Update không được chứa hai bước cần nhớ trạng thái liên tiếp, ví
dụ \enquote{loại trùng rồi đếm}~\cite{sparkguide}. Vì vậy WikiPulse tách thành nhiều truy vấn: một truy vấn loại trùng và
ghi ra đĩa, các truy vấn khác đọc lại kết quả đó để đếm (Mục~\ref{sec:spark}).

\subsection{Loại bản trùng}
Mỗi sự kiện có một định danh duy nhất. Spark nhớ các định danh đã thấy và bỏ sự kiện nào trùng. Không thể nhớ mãi, nên
Spark chỉ nhớ trong phạm vi watermark: bản trùng đến cách bản gốc quá 10 phút thì không còn bị phát hiện. Đây là giới hạn
có chủ đích, được ghi rõ trong contract của nhóm.

\section{Kiến trúc dữ liệu ba tầng}
Thay vì chỉ lưu kết quả cuối, dữ liệu được lưu theo ba tầng~\cite{medallion}, mỗi tầng sạch hơn tầng trước
(Hình~\ref{fig:medallion}). Lưu tầng thô giúp kiểm tra lại khi nghi có lỗi và tính lại khi đổi công thức mà không cần nguồn
gửi lại dữ liệu.

\begin{figure}[H]
\centering
\begin{tikzpicture}[font=\footnotesize\sffamily]
  \node[draw=wpamber,fill=lightamber,rounded corners=3pt,text width=38mm,align=left,inner sep=5pt] (b) at (0,0)
    {\textbf{Bronze -- thô}\\[2pt] Nguyên văn mọi bản ghi nhận được, kể cả bản hỏng (tách riêng).\\[2pt]
     \textcolor{wpgray}{Còn tên người dùng, IP.}};
  \node[draw=wpgray,fill=lightgray,rounded corners=3pt,text width=38mm,align=left,inner sep=5pt] (s) at (5.2,0)
    {\textbf{Silver -- đã làm sạch}\\[2pt] Đã kiểm tra, chuẩn hóa kiểu dữ liệu, giờ UTC, bỏ trùng, bỏ canary.\\[2pt]
     \textcolor{wpgray}{Người dùng chỉ còn mã băm.}};
  \node[draw=wpamber!60!black,fill=wpamber!25,rounded corners=3pt,text width=38mm,align=left,inner sep=5pt] (g) at (10.4,0)
    {\textbf{Gold -- tổng hợp}\\[2pt] Số đếm theo wiki và cửa sổ, chỉ ghi khi đã chốt.\\[2pt]
     \textcolor{wpgray}{Dùng cho phân tích, báo cáo.}};
  \draw[arr] (b) -- node[above,lab]{kiểm tra} (s);
  \draw[arr] (s) -- node[above,lab]{tổng hợp} (g);
\end{tikzpicture}
\caption{Ba tầng dữ liệu. Đi từ trái sang phải, dữ liệu sạch dần và ít thông tin nhạy cảm hơn.}\label{fig:medallion}
\end{figure}

Ba tầng được lưu dạng \textbf{Parquet}, định dạng theo cột, nén tốt, đọc nhanh khi chỉ cần vài cột. File được chia thư
mục theo ngày và giờ, nên truy vấn một khoảng thời gian chỉ phải đọc đúng các thư mục đó.

\section{Apache Cassandra: phục vụ kết quả}
Dashboard cần đọc số liệu nhanh, liên tục. Parquet hợp cho phân tích theo lô nhưng không hợp để trả lời hàng nghìn câu
hỏi nhỏ mỗi phút. Kết quả tạm tính vì vậy được ghi vào Cassandra~\cite{lakshman2010}.

Khác CSDL quan hệ, trong Cassandra \textbf{bảng được thiết kế theo câu truy vấn}, không theo thực thể~\cite{cassandracql}.
Khóa chính có hai phần (Hình~\ref{fig:cassandra}):
\begin{itemize}
  \item \textbf{Partition key} quyết định dữ liệu nằm cùng nhau ở đâu. WikiPulse dùng wiki: mọi cửa sổ của enwiki nằm
        chung một partition.
  \item \textbf{Clustering key} quyết định thứ tự trong partition. WikiPulse dùng thời điểm bắt đầu cửa sổ, sắp giảm dần,
        nên các phút mới nhất luôn nằm đầu.
\end{itemize}

\begin{figure}[H]
\centering
\begin{tikzpicture}[font=\footnotesize\sffamily]
  \foreach \p/\x/\rows in {enwiki/0/{00:12 | 1 sự kiện, 00:00 | 5 sự kiện},
                           viwiki/6.6/{00:00 | 1 sự kiện}} {
    \node[font=\footnotesize\sffamily\bfseries,anchor=west] at (\x,0.55) {partition \enquote{\p}};
    \foreach \r [count=\i] in \rows {
      \node[draw=wpblue,fill=lightblue,minimum width=50mm,minimum height=6mm,anchor=west] at (\x,-\i*0.65+0.1) {\r};
    }
  }
  \draw[arr] (5.3,0.05) -- node[right,lab,align=left]{mới\\cũ} (5.3,-1.3);
  \node[draw=wpgreen,fill=lightgreen,rounded corners=3pt,text width=100mm,align=left,inner sep=4pt] at (5.8,-2.3)
    {\enquote{5 phút gần nhất của enwiki} = vào đúng partition enwiki, lấy 5 dòng đầu. Không phải quét bảng.};
\end{tikzpicture}
\caption{Dữ liệu hoạt động trong Cassandra. Mỗi wiki là một partition; trong partition, các cửa sổ xếp mới trước.}\label{fig:cassandra}
\end{figure}

Hai tính chất của Cassandra quyết định cách ghi:
\begin{itemize}
  \item Ghi một dòng có khóa đã tồn tại thì \textbf{ghi đè} (upsert), không sinh dòng mới. WikiPulse ghi \emph{tổng số}
        của cửa sổ chứ không ghi \enquote{cộng thêm 1}, nên khi Spark ghi lại một batch sau sự cố, kết quả vẫn đúng, không
        bị đếm đôi.
  \item Truy vấn không chỉ rõ partition buộc Cassandra quét mọi node. Dự án cấm kiểu truy vấn này (\enquote{allow
        filtering}); mỗi câu hỏi của API phải có bảng thiết kế riêng cho nó.
\end{itemize}

\section{FastAPI và Streamlit}
\textbf{FastAPI} là framework Python để viết API web; kiểu dữ liệu khai báo một lần được dùng để kiểm tra đầu vào, định
dạng đầu ra và sinh tài liệu API tự động (OpenAPI). \textbf{Streamlit} dựng giao diện dashboard bằng Python thuần, hợp với
nhóm không có người chuyên frontend.

\section{Công nghệ sử dụng}
\begin{table}[H]
\centering\small
\caption{Công nghệ sử dụng và vai trò}\label{tab:congnghe}
\begin{tabularx}{\linewidth}{@{}llX@{}}
\toprule
\textbf{Thành phần} & \textbf{Lựa chọn} & \textbf{Vai trò} \\ \midrule
Hàng đợi & Apache Kafka 3.9 & Giữ sự kiện, tách nguồn khỏi xử lý, cho phép đọc lại \\
Xử lý luồng & Spark 3.5 (PySpark) & Kiểm tra, chuẩn hóa, loại trùng, đếm theo cửa sổ \\
Lưu trữ phân tích & Parquet (cục bộ, sau là HDFS) & Ba tầng Bronze, Silver, Gold \\
Lưu trữ phục vụ & Apache Cassandra 4.1 & Số tạm tính cho API \\
API & FastAPI & Cung cấp số liệu cho dashboard \\
Giao diện & Streamlit & Dashboard bốn trang \\
Chạy & Docker Compose & Dựng toàn bộ hệ thống bằng một lệnh, chạy trên ARM64 và x86\_64 \\
Ngôn ngữ, kiểm thử & Python 3.11, pytest & Toàn bộ mã nguồn và kiểm thử \\
\bottomrule
\end{tabularx}
\end{table}
\dukiensua{Thêm HDFS khi có profile đầy đủ (tuần 2--3).}

% ============================================================================
% CHƯƠNG 3: LIÊN QUAN
% ============================================================================
\chapter{Các công cụ và nghiên cứu liên quan}\label{chap:lienquan}
\section{Công cụ tuần tra của cộng đồng Wikipedia}\chualam
\section{Mô hình đánh giá chỉnh sửa của Wikimedia}\chualam
\section{Thống kê lượt xem và xu hướng trang}\chualam
\section{Khoảng trống mà WikiPulse nhắm tới}\chualam

% ============================================================================
% CHƯƠNG 4: PHÂN TÍCH, THIẾT KẾ VÀ CÀI ĐẶT
% ============================================================================
\chapter{Phân tích, thiết kế và cài đặt}\label{chap:thietke}

\section{Phân tích yêu cầu}
\subsection{Tác nhân}
\begin{itemize}
  \item \textbf{Người xem:} người tuần tra hoặc người quan tâm đến Wikipedia, mở dashboard để xem hoạt động, trang xu
        hướng và các chỉnh sửa nên xem trước. Không cần đăng nhập (ngoài phạm vi MVP).
  \item \textbf{Người vận hành:} thành viên nhóm khởi động, dừng hệ thống, phát lại dữ liệu mẫu, chuyển sang nguồn thật
        và kiểm tra trạng thái. Người vận hành làm được mọi việc của người xem.
  \item \textbf{Wikimedia EventStreams (hệ thống ngoài):} phát luồng sự kiện thay đổi ở chế độ LIVE.
\end{itemize}

\subsection{Biểu đồ use case}
\begin{figure}[H]
\centering
\begin{tikzpicture}[font=\small\sffamily]
  \draw[thick,rounded corners=3pt] (-3.4,1.0) rectangle (3.4,-9.9);
  \node[font=\sffamily\bfseries] at (0,0.6) {Hệ thống WikiPulse};
  \node[uc] (u1) at (0,-0.2) {Xem hoạt động theo wiki};
  \node[uc,dashed] (u2) at (0,-1.2) {Xem trang xu hướng};
  \node[uc,dashed] (u3) at (0,-2.2) {Xem chỉnh sửa cần ưu tiên};
  \node[uc,dashed] (u4) at (0,-3.2) {Xem chi tiết một sự kiện};
  \node[uc] (u5) at (0,-4.2) {Xem trạng thái hệ thống};
  \node[uc] (u6) at (0,-5.5) {Khởi động, dừng hệ thống};
  \node[uc] (u7) at (0,-6.5) {Phát lại dữ liệu mẫu};
  \node[uc] (u8) at (0,-7.5) {Chạy kiểm thử end-to-end};
  \node[uc,dashed] (u9) at (0,-8.9) {Thu thập luồng thật (LIVE)};
  \pic at (-6,-1.9) {actor}; \node (nx) at (-6,-1.9) {}; \node[left] at (-6.35,-1.9) {Người xem};
  \pic at (-6,-6.4) {actor}; \node (vh) at (-6,-6.4) {}; \node[left,align=right] at (-6.35,-6.4) {Người\\vận hành};
  \pic at (6,-8.9) {actor}; \node (wm) at (6,-8.9) {}; \node[below,align=center] at (6,-9.4) {Wikimedia\\EventStreams};
  \foreach \u in {1,...,5} \draw (nx) -- (u\u.west);
  \foreach \u in {6,7,8,9} \draw (vh) -- (u\u.west);
  \draw (wm) -- (u9.east);
  \draw[-{Triangle[open,length=3mm,width=3mm]},thick] (-6,-5.85) -- (-6,-2.45);
\end{tikzpicture}
\caption{Biểu đồ use case tổng quát. Đường nối là quan hệ giữa tác nhân và use case; mũi tên tam giác rỗng là quan hệ
khái quát hóa (người vận hành làm được mọi việc của người xem). Use case viền đứt nét là phần chưa làm.}\label{fig:usecase}
\end{figure}
\dukiensua{Viền đứt nét đổi thành liền khi use case tương ứng xong (tuần 2--3).}

\subsection{Đặc tả các use case chính}
\begin{table}[H]
\centering\small
\caption{Đặc tả use case \enquote{Xem hoạt động theo wiki}}\label{tab:uc-hoatdong}
\begin{tabularx}{\linewidth}{@{}>{\bfseries}p{3.2cm}X@{}}
\toprule
Tác nhân & Người xem \\
Tiền điều kiện & Hệ thống đang chạy \\
Luồng chính & 1. Người xem mở trang Overview, chọn wiki (mặc định enwiki).\newline
2. Dashboard gọi API lấy các cửa sổ 1 phút mới nhất của wiki đó.\newline
3. API đọc đúng partition của wiki trong Cassandra và trả kết quả, gắn nhãn \enquote{tạm tính}.\newline
4. Dashboard hiện năm chỉ số (sự kiện, chỉnh sửa, trang mới, sửa ẩn danh, sửa bởi bot) và biểu đồ theo thời gian;
tự làm mới sau 15 giây.\newline
5. Người xem có thể chọn một cửa sổ cụ thể để xem riêng. \\
Luồng thay thế & 3a. Chưa có dữ liệu: hiện thông báo trống, không hiện số giả.\newline
3b. Cassandra không phản hồi: API trả lỗi \enquote{phụ thuộc chưa sẵn sàng}; dashboard báo dịch vụ chưa sẵn sàng.\newline
1a. Wiki không nằm trong danh sách cho phép: API từ chối yêu cầu. \\
Hậu điều kiện & Không thay đổi dữ liệu \\
\bottomrule
\end{tabularx}
\end{table}

\begin{table}[H]
\centering\small
\caption{Đặc tả use case \enquote{Phát lại dữ liệu mẫu}}\label{tab:uc-replay}
\begin{tabularx}{\linewidth}{@{}>{\bfseries}p{3.2cm}X@{}}
\toprule
Tác nhân & Người vận hành \\
Tiền điều kiện & Kafka đã sẵn sàng và đã có topic \\
Luồng chính & 1. Người vận hành chạy lệnh phát lại.\newline
2. Hệ thống đọc từng dòng của file dữ liệu mẫu, giữ nguyên định danh và thời điểm của sự kiện, thêm thời điểm nhận.\newline
3. Hệ thống gửi từng sự kiện vào Kafka với khóa wiki và trang, chờ Kafka xác nhận.\newline
4. Spark xử lý như với dữ liệu thật; sau khoảng 10--30 giây số liệu xuất hiện trên dashboard. \\
Luồng thay thế & 2a. Dòng không phải JSON hợp lệ: vẫn gửi nguyên văn để kiểm tra đường xử lý bản ghi hỏng.\newline
1a. Người vận hành chọn tốc độ phát nhanh hơn: chỉ đổi nhịp gửi, không đổi thời điểm sự kiện. \\
Hậu điều kiện & Kafka có thêm các sự kiện mẫu; phát lại nhiều lần không làm tăng số đếm trong phạm vi watermark \\
\bottomrule
\end{tabularx}
\end{table}

\subsubsection{Đặc tả use case \enquote{Xem trang xu hướng}}\chualam
\subsubsection{Đặc tả use case \enquote{Xem chỉnh sửa cần ưu tiên}}\chualam
\subsubsection{Đặc tả use case \enquote{Xem chi tiết một sự kiện}}\chualam
\subsubsection{Đặc tả use case \enquote{Thu thập luồng thật (LIVE)}}\chualam

\subsection{Yêu cầu chức năng}
\begin{table}[H]
\centering\small
\caption{Yêu cầu chức năng và trạng thái}\label{tab:ycchucnang}
\begin{tabularx}{\linewidth}{@{}lXl@{}}
\toprule
\textbf{Mã} & \textbf{Yêu cầu} & \textbf{Trạng thái} \\ \midrule
F1 & Phát lại dữ liệu mẫu vào Kafka, giữ nguyên định danh và thời điểm sự kiện & Đã làm \\
F2 & Nhận luồng thật từ Wikimedia, tự nối lại khi mất kết nối & Tuần 2 \\
F3 & Kiểm tra, tách riêng bản ghi hỏng, chuẩn hóa, loại trùng & Đã làm \\
F4 & Đếm hoạt động theo wiki, cửa sổ 1 phút; có số tạm tính và số đã chốt & Đã làm \\
F5 & Trang xu hướng: cửa sổ 5 phút trượt mỗi 1 phút, top-N theo wiki & Tuần 2 \\
F6 & Ưu tiên xem xét chỉnh sửa theo ba dấu hiệu: xóa nhiều nội dung, hoàn tác qua lại, sửa ẩn danh & Tuần 2 \\
F7 & Tra cứu chi tiết một sự kiện & Tuần 2--3 \\
F8 & Dashboard 4 trang: Overview, Trending, Suspicious Edits, System Status & Overview đã làm \\
F9 & Lưu ba tầng trên HDFS & Tuần 2--3 \\
\bottomrule
\end{tabularx}
\end{table}
\dukiensua{Cập nhật trạng thái theo từng checkpoint.}

\subsection{Yêu cầu phi chức năng}
\begin{itemize}
  \item \textbf{Độ trễ:} số tạm tính hiện trên API trong 30--60 giây ở tải demo (mục tiêu, chưa đo). Số đã chốt chậm hơn
        theo watermark 10 phút.
  \item \textbf{Đúng khi chạy lại:} retry hay khởi động lại không được làm số bị đếm đôi.
  \item \textbf{Riêng tư:} ngoài tầng thô, không nơi nào chứa tên người dùng hay IP. Người dùng được thay bằng mã băm
        có \enquote{muối}; đây là bút danh hóa, không phải ẩn danh hóa hoàn toàn.
  \item \textbf{Tái hiện được:} mọi kết quả kiểm thử chạy lại được bằng lệnh có sẵn.
  \item \textbf{Tài nguyên:} chạy được trên máy 16 GB RAM, cấp cho Docker 6--8 GB.
\end{itemize}

\section{Kiến trúc tổng quan}
\begin{figure}[H]
\centering
\begin{tikzpicture}[scale=0.86, transform shape]
  \node[abox=24mm] (src) at (0,0) {Wikimedia\\EventStreams};
  \node[box=24mm] (rep) at (0,-2.2) {Replay\\(dữ liệu mẫu)};
  \node[box=24mm,dashed] (col) at (3.6,0) {Collector LIVE};
  \node[db] (kafka) at (3.6,-2.2) {Kafka\\sự kiện thô};
  \node[box=30mm] (spark) at (7.6,-2.2) {Spark Structured\\Streaming};
  \node[db] (bronze) at (12.2,0.2) {Bronze\\thô, hỏng};
  \node[db] (silver) at (12.2,-2.2) {Silver\\đã làm sạch};
  \node[db] (gold) at (12.2,-4.6) {Gold\\đã chốt};
  \node[db] (cass) at (7.6,-4.9) {Cassandra\\tạm tính};
  \node[gbox=24mm] (api) at (3.6,-4.9) {FastAPI};
  \node[gbox=24mm] (dash) at (0,-4.9) {Dashboard\\Streamlit};
  \draw[arr,dashed] (src) -- (col);
  \draw[arr,dashed] (col) -- (kafka);
  \draw[arr] (rep) -- (kafka);
  \draw[arr] (kafka) -- (spark);
  \draw[arr] (spark) -- (bronze.west);
  \draw[arr] (spark) -- (silver);
  \draw[arr] (spark) -- node[below,lab]{append} (gold.west);
  \draw[arr] (spark) -- node[right,lab]{update} (cass);
  \draw[arr] (cass) -- (api);
  \draw[arr] (api) -- (dash);
\end{tikzpicture}
\caption{Sơ đồ kiến trúc tổng quan. Nét đứt là phần chưa làm (collector LIVE, tuần 2). Màu: xanh dương là thành phần
xử lý, xanh lá là đầu ra cho người dùng, vàng là nguồn bên ngoài.}\label{fig:kientruc}
\end{figure}

Kafka tách phần nhận khỏi phần xử lý; Spark là nơi duy nhất biến đổi dữ liệu; API và dashboard chỉ đọc kết quả. Dashboard
không đọc Cassandra trực tiếp mà luôn đi qua API, để API là điểm duy nhất áp quy tắc riêng tư và giới hạn truy vấn.

\subsection{Các thành phần và người phụ trách}
\begin{table}[H]
\centering\small
\caption{Các thành phần (vị trí mã nguồn ở Phụ lục~\ref{app:ma})}\label{tab:module}
\begin{tabularx}{\linewidth}{@{}lXl@{}}
\toprule
\textbf{Thành phần} & \textbf{Nhiệm vụ} & \textbf{Phụ trách} \\ \midrule
Lõi dùng chung & Đọc và kiểm tra cấu hình, ghi log dạng JSON, lớp truy cập Cassandra & Trọng, Nhân \\
Thu thập & Phát lại dữ liệu mẫu, thu thập luồng thật & Thành \\
Xử lý & Tiến trình Spark: các truy vấn, phép biến đổi, giám sát truy vấn & Linh \\
Ghi kết quả & Đưa số đếm từ Spark vào Cassandra & Nhân \\
Phục vụ & Lược đồ Cassandra, API & Nhân \\
Giao diện & Dashboard và lớp gọi API & Mạnh \\
Tích hợp & Docker Compose, tạo topic, dữ liệu mẫu, kiểm thử end-to-end & Trọng \\
\bottomrule
\end{tabularx}
\end{table}

\subsection{Triển khai}
\begin{figure}[H]
\centering
\begin{tikzpicture}[scale=0.8, transform shape, font=\footnotesize\sffamily]
  \node[udnode,text width=22mm] (kafka) at (0,0) {«container»\\\textbf{kafka}};
  \node[udnode,text width=22mm] (spark) at (4,0) {«container»\\\textbf{spark}};
  \node[udnode,text width=22mm] (cass) at (8,0) {«container»\\\textbf{cassandra}};
  \node[udnode,text width=22mm] (api) at (8,-2.8) {«container»\\\textbf{api}\\cổng 18000};
  \node[udnode,text width=22mm] (dash) at (4,-2.8) {«container»\\\textbf{dashboard}\\cổng 8501};
  \node[udnode,text width=22mm] (rep) at (0,-2.8) {«container»\\\textbf{replay}\\chạy khi cần};
  \begin{scope}[on background layer]
    \node[udnode,fill=white,fit=(kafka)(cass)(rep)(api),inner xsep=6mm,inner ysep=6mm,yshift=4mm,
          label={[font=\footnotesize\sffamily,anchor=north]north:«device» \textbf{Máy phát triển} (Docker Compose)}] (host) {};
  \end{scope}
  \node[udnode,text width=22mm] (browser) at (4,-6.2) {«device»\\\textbf{Trình duyệt}};
  \draw[thick] (rep) -- node[left,guard]{Kafka} (kafka);
  \draw[thick] (kafka) -- node[above,guard]{Kafka} (spark);
  \draw[thick] (spark) -- node[above,guard]{CQL} (cass);
  \draw[thick] (api) -- node[right,guard]{CQL} (cass);
  \draw[thick] (dash) -- node[above,guard]{HTTP} (api);
  \draw[thick] (browser) -- node[right,guard]{HTTP} (dash);
\end{tikzpicture}
\caption{Biểu đồ triển khai hiện tại. Các container gọi nhau bằng tên dịch vụ trong mạng Docker; chỉ API, dashboard và
Kafka mở cổng ra máy chủ, và chỉ cho chính máy đó truy cập. Hai container khởi tạo (tạo topic, tạo lược đồ Cassandra)
chạy một lần trước các dịch vụ chính.}\label{fig:trienkhai}
\end{figure}
\dukiensua{Thêm NameNode/DataNode khi có profile HDFS (W2-TRONG-01); Kubernetes ở tuần 3.}

\section{Đường đi của một sự kiện}
Hình~\ref{fig:hoatdong} cho thấy một sự kiện đi qua những bước kiểm tra nào trước khi được đếm. Mỗi nhánh rẽ đều để
lại dấu vết: không có sự kiện nào \enquote{biến mất} mà không được lưu ở tầng thô.

\begin{figure}[H]
\centering
\begin{tikzpicture}[font=\footnotesize\sffamily, node distance=5mm]
  \node[start] (s) at (0,0) {};
  \node[act=50mm] (a1) at (0,-0.9) {Nhận sự kiện từ Kafka};
  \node[act=50mm] (a2) at (0,-2.0) {Lưu nguyên văn vào Bronze};
  \node[dnode] (d1) at (0,-3.2) {};
  \node[act=40mm,fill=lightred] (x1) at (5.6,-3.2) {Lưu vào vùng bản ghi hỏng\\kèm mã lỗi};
  \node[dnode] (d2) at (0,-4.8) {};
  \node[act=40mm,fill=lightgray] (x2) at (5.6,-4.8) {Bỏ qua\\(chỉ còn ở Bronze)};
  \node[act=62mm] (a3) at (0,-6.3) {Chuẩn hóa: giờ UTC, băm người dùng,\\tính chênh lệch byte, gán định danh};
  \node[dnode] (d3) at (0,-7.8) {};
  \node[act=40mm,fill=lightgray] (x3) at (5.6,-7.8) {Bỏ bản trùng};
  \node[act=50mm] (a4) at (0,-9.0) {Ghi vào Silver};
  \node[act=50mm] (a5) at (0,-10.1) {Cộng vào cửa sổ 1 phút của wiki};
  \node[act=40mm,fill=lightgreen] (o1) at (-5.0,-11.2) {Ghi số tạm tính\\vào Cassandra (mỗi batch)};
  \node[act=40mm,fill=lightamber] (o2) at (5.0,-11.2) {Ghi số đã chốt vào Gold\\(khi qua watermark)};
  \node[stop] (e) at (0,-12.4) {};
  \draw[flow] (s) -- (a1); \draw[flow] (a1) -- (a2); \draw[flow] (a2) -- (d1);
  \draw[flow] (d1) -- node[above,guard]{[hỏng]} (x1);
  \draw[flow] (d1) -- node[left,guard]{[hợp lệ]} (d2);
  \draw[flow] (d2) -- node[above,guard,align=center]{[wiki ngoài danh sách\\hoặc canary]} (x2);
  \draw[flow] (d2) -- node[left,guard]{[còn lại]} (a3);
  \draw[flow] (a3) -- (d3);
  \draw[flow] (d3) -- node[above,guard]{[đã thấy định danh]} (x3);
  \draw[flow] (d3) -- node[left,guard]{[mới]} (a4);
  \draw[flow] (a4) -- (a5);
  \draw[flow] (a5.west) -| (o1.north);
  \draw[flow] (a5.east) -| (o2.north);
  \draw[flow] (o1.south) |- (e.west);
  \draw[flow] (o2.south) |- (e.east);
\end{tikzpicture}
\caption{Biểu đồ hoạt động: đường đi của một sự kiện. Hình thoi là nút rẽ nhánh, điều kiện ghi trong ngoặc vuông.
Bốn loại lỗi được ghi mã riêng: không phải JSON hợp lệ, thiếu thời điểm, thiếu wiki, thiếu trang.}\label{fig:hoatdong}
\end{figure}

\section{Biểu đồ tuần tự: người xem mở trang Overview}\chualam
\dukiensua{Vẽ sau khi W1-MANH-01 xong, lúc luồng dashboard $\rightarrow$ API $\rightarrow$ Cassandra đã ổn định.}

\section{Cấu hình và tách lần chạy}
Mọi tham số (chế độ LIVE hay REPLAY, danh sách wiki, độ trễ watermark, địa chỉ các dịch vụ\ldots) được đọc từ biến môi
trường và kiểm tra ngay khi khởi động, trước khi mở kết nối tốn tài nguyên. Mỗi lần chạy có một \textbf{mã lần chạy};
dữ liệu và checkpoint của lần chạy nằm trong thư mục riêng mang mã đó, và mỗi truy vấn Spark có checkpoint riêng. Khi đổi
schema hay nguồn dữ liệu, nhóm tạo mã lần chạy và keyspace Cassandra mới chứ không xóa checkpoint rồi ghi tiếp vào kết quả
cũ. Kiểm thử end-to-end dùng mã lần chạy và vùng lưu trữ riêng nên không đụng dữ liệu đang dùng để phát triển.

\section{Thu thập dữ liệu}
\subsection{Kafka}
Hệ thống dùng ba topic: sự kiện thô (3 partition), bản ghi không xử lý được (dead-letter queue) và cảnh báo của đường
ống; một broker, mỗi partition một bản sao. Sự kiện giữ nguyên JSON của Wikimedia, chỉ thêm thời điểm nhận và phiên bản
bộ thu thập. Khóa là wiki và mã trang (thiếu mã trang thì dùng tên trang).

\subsection{Phát lại dữ liệu mẫu}
Bộ phát lại đọc file dữ liệu mẫu và gửi từng dòng vào Kafka, giữ nguyên định danh và thời điểm sự kiện; tham số tốc độ
chỉ đổi nhịp gửi. Dữ liệu mẫu ngày 01/01/2026 có 11 dòng, được chọn để thử đủ các nhánh ở Hình~\ref{fig:hoatdong}: 5
chỉnh sửa enwiki, 1 trang mới viwiki, 1 bản trùng, 2 bản hỏng, 1 sự kiện canary, và 1 sự kiện lúc 00:12:10 để đẩy
watermark qua cửa sổ 00:00 (Hình~\ref{fig:watermark}).

\subsection{Thu thập luồng thật}\chualam

\section{Xử lý luồng}\label{sec:spark}
\subsection{Các truy vấn}
Một tiến trình Spark chạy năm truy vấn độc lập, mỗi truy vấn một checkpoint (Bảng~\ref{tab:truyvan}). Bước loại trùng
ghi kết quả ra Silver; hai truy vấn đếm đọc lại Silver. Nhờ vậy không truy vấn nào vừa loại trùng vừa đếm ở chế độ Update
(giới hạn đã nêu ở Chương~\ref{chap:nen}).

\begin{table}[H]
\centering\small
\caption{Các truy vấn streaming}\label{tab:truyvan}
\begin{tabularx}{\linewidth}{@{}lXl@{}}
\toprule
\textbf{Truy vấn} & \textbf{Việc làm} & \textbf{Ghi vào (chế độ)} \\ \midrule
Bronze thô & Đọc Kafka, giữ nguyên bản ghi và vị trí trong Kafka & Parquet (append) \\
Bronze hỏng & Đọc Kafka, lọc bản ghi hỏng, gắn mã lỗi & Parquet (append) \\
Silver & Đọc Kafka, lọc, chuẩn hóa, loại trùng & Parquet (append) \\
Hoạt động tạm tính & Đọc Silver, đếm theo wiki và phút & Cassandra (update) \\
Hoạt động đã chốt & Đọc Silver, đếm như trên & Parquet Gold (append) \\
\bottomrule
\end{tabularx}
\end{table}

\subsection{Định danh và loại trùng}
Định danh của sự kiện lấy từ nguồn. Nếu nguồn không có, hệ thống tự tạo bằng cách băm các trường không đổi của sự kiện
(wiki, trang, phiên bản, thời điểm, loại); không bao giờ dùng thời điểm nhận hay xử lý, vì như vậy cùng một sự kiện gửi
hai lần sẽ có hai định danh khác nhau. Mọi phép biến đổi dùng hàm có sẵn của Spark, không viết hàm Python riêng, để Spark
tối ưu được và không phải chuyển dữ liệu qua lại giữa Java và Python.

\subsection{Đếm hoạt động}
Với mỗi wiki và mỗi phút, hệ thống đếm: tổng sự kiện hợp lệ, số chỉnh sửa, số trang mới, và trong số chỉnh sửa thì bao
nhiêu là ẩn danh, bao nhiêu là của bot.

\subsection{Đặc trưng xu hướng}\chualam
\subsection{Chấm điểm ưu tiên xem xét}\chualam

\section{Lưu trữ}
\subsection{Tầng Parquet}
Bronze và Silver chia thư mục theo ngày và giờ của sự kiện; Gold chia theo ngày, giờ và wiki. Danh sách đầy đủ các trường
của Silver nằm trong tài liệu contract của nhóm.
\dukiensua{Chuyển sang HDFS ở tuần 2--3 (W2-TRONG-01, W3-LINH-01); thêm kết quả đọc lại sau khởi động lại.}

\subsection{Lược đồ Cassandra}\label{sec:cassandra}
Bảng hoạt động theo cửa sổ (Hình~\ref{fig:bangcass}) phục vụ đúng hai câu hỏi của API: \enquote{các phút mới nhất của
một wiki} và \enquote{một phút cụ thể của một wiki}. Cách chọn khóa đã giải thích ở Hình~\ref{fig:cassandra}.

\begin{figure}[H]
\centering
\begin{tikzpicture}[scale=1.5, transform shape]
\node[erd] (t) at (0,0) {\textbf{wiki\_activity\_by\_window}\nodepart{two}\begin{tabular}{@{}l@{\;}l@{\;}l@{\;}l@{}}
PK & \textbf{wiki} & {\color{wpgray}text} & khóa partition \\
CK$\downarrow$ & \textbf{window\_start} & {\color{wpgray}timestamp} & sắp mới trước \\
 & window\_end & {\color{wpgray}timestamp} & \\
 & event\_count & {\color{wpgray}bigint} & tổng sự kiện \\
 & edit\_count & {\color{wpgray}bigint} & chỉnh sửa \\
 & new\_page\_count & {\color{wpgray}bigint} & trang mới \\
 & anonymous\_edit\_count & {\color{wpgray}bigint} & sửa ẩn danh \\
 & bot\_edit\_count & {\color{wpgray}bigint} & sửa bởi bot \\
 & updated\_at & {\color{wpgray}timestamp} & lần ghi gần nhất \\
\end{tabular}};
\end{tikzpicture}
\caption{Bảng hoạt động theo cửa sổ. PK: khóa partition; CK$\downarrow$: khóa sắp xếp, giảm dần.}\label{fig:bangcass}
\end{figure}

\begin{itemize}
  \item Spark ghi \textbf{tổng số} của cửa sổ, không ghi phần cộng thêm, nên ghi lại sau sự cố không làm đếm đôi.
  \item Mỗi wiki thêm 1.440 dòng mỗi ngày, khoảng 525 nghìn dòng mỗi năm: đủ cho demo. Trước khi chạy LIVE dài ngày cần
        đặt thời hạn tự xóa (TTL) hoặc chia partition theo ngày.
  \item Trước khi ghi, mỗi dòng được kiểm tra: tên wiki hợp lệ, cửa sổ kết thúc sau khi bắt đầu, số đếm nhất quán (chỉnh
        sửa cộng trang mới không vượt tổng; ẩn danh, bot không vượt số chỉnh sửa). Dòng sai làm batch thất bại để Spark
        chạy lại, thay vì ghi số sai.
\end{itemize}
\dukiensua{Thêm bảng xu hướng, chỉnh sửa đáng ngờ, chi tiết sự kiện ở tuần 2 (W2-NHAN-01); TTL ở tuần 3 (W3-NHAN-01).}

\subsection{Ghi kết quả từ Spark sang Cassandra}
Ở mỗi batch, tiến trình chính của Spark duyệt lần lượt từng dòng kết quả (không nạp cả batch vào bộ nhớ) và ghi vào
Cassandra. Nếu ghi lỗi, truy vấn dừng và Spark chạy lại batch đó từ checkpoint. Parquet và Cassandra không có giao dịch
chung, nên hệ thống không cam kết \enquote{đúng một lần} trên cả hai nơi; tính đúng dựa vào việc ghi đè theo khóa ổn định.

\section{API}
\begin{table}[H]
\centering\small
\caption{Các endpoint}\label{tab:api}
\begin{tabularx}{\linewidth}{@{}lXl@{}}
\toprule
\textbf{Endpoint} & \textbf{Mô tả} & \textbf{Trạng thái} \\ \midrule
\code{/health} & Kiểm tra Cassandra thật sự trả lời được & Đã làm \\
\code{/api/v1/activity} & Hoạt động theo cửa sổ; chọn wiki, một cửa sổ cụ thể, số dòng tối đa 100 & Đã làm \\
\code{/api/v1/trends} & Top-N trang xu hướng & Tuần 2 \\
\code{/api/v1/suspicious-edits} & Chỉnh sửa cần ưu tiên xem xét & Tuần 2 \\
\code{/api/v1/events/\{id\}} & Chi tiết một sự kiện & Tuần 2--3 \\
\bottomrule
\end{tabularx}
\end{table}

Mọi phản hồi dữ liệu kèm thời điểm tạo, nguồn, chế độ chạy và nhãn \enquote{tạm tính}; mọi thời điểm theo UTC. Không có
dữ liệu thì trả danh sách rỗng, không phải lỗi. Tham số sai (wiki ngoài danh sách, cửa sổ thiếu múi giờ, số dòng ngoài
1--100) bị từ chối. Cassandra không trả lời thì API báo \enquote{phụ thuộc chưa sẵn sàng}; API vẫn khởi động được khi
Cassandra chưa lên và tự hồi phục khi Cassandra sẵn sàng. Endpoint chưa làm trả mã \enquote{chưa triển khai} kèm người
phụ trách và tuần, để dashboard hiện đúng trạng thái thay vì dữ liệu giả.
\dukiensua{Thống nhất định dạng lỗi cho mọi endpoint và thêm ba endpoint ở tuần 2.}

\section{Dashboard}
Bốn trang: Overview (đã đọc dữ liệu thật: chọn wiki và cửa sổ, tự làm mới sau 15 giây), Trending, Suspicious Edits,
System Status. Các trang chưa có dữ liệu hiện rõ là chưa triển khai, kèm người phụ trách.
\dukiensua{Mạnh bổ sung ảnh chụp các trang (W1-MANH-01, W2-MANH-01).}

% ============================================================================
% CHƯƠNG 5: KIỂM THỬ VÀ ĐÁNH GIÁ
% ============================================================================
\chapter{Kiểm thử và đánh giá}\label{chap:danhgia}

\section{Chiến lược kiểm thử}
\begin{table}[H]
\centering\small
\caption{Các tầng kiểm thử}\label{tab:kiemthu}
\begin{tabularx}{\linewidth}{@{}lXl@{}}
\toprule
\textbf{Tầng} & \textbf{Nội dung} & \textbf{Khi chạy} \\ \midrule
Nhẹ & Cấu hình, API, lớp truy cập Cassandra, phát lại, dashboard; Cassandra được giả lập ở tầng driver nhưng giữ đúng cách ghi đè theo khóa & Mỗi lần sửa code, vài giây \\
Spark & Các phép biến đổi và cách điều phối truy vấn, chạy trong container có Java & Khi sửa phần xử lý \\
Tích hợp & Lớp truy cập trên Cassandra thật, keyspace ngẫu nhiên, tự xóa sau khi chạy & Khi sửa lược đồ, truy vấn \\
End-to-end & Dựng hệ thống riêng, phát dữ liệu mẫu hai pha, kiểm tra cả ba tầng, API và giao diện & Trước khi nghiệm thu \\
\bottomrule
\end{tabularx}
\end{table}
Lệnh chạy từng tầng ở Phụ lục~\ref{app:ma}.

\section{Kết quả CP0}
Dữ liệu mẫu ngày 01/01/2026 gồm 11 bản ghi. Kết quả chạy thật (Bảng~\ref{tab:cp0}) khớp bộ kết quả kỳ vọng; bộ này
được viết tay độc lập, không sinh ra từ chính code xử lý.

\begin{table}[H]
\centering\small
\caption{Kết quả CP0 trên dữ liệu mẫu, cửa sổ 2026-01-01 00:00 UTC}\label{tab:cp0}
\begin{tabularx}{\linewidth}{@{}Xr@{}}
\toprule
\textbf{Chỉ tiêu} & \textbf{Kết quả} \\ \midrule
Bản ghi Kafka / Bronze raw & 11 / 11 \\
Bronze invalid & 2 \\
Silver (sau loại trùng và canary) & 7 \\
Gold cửa sổ 00:00 (đã chốt) & 2 dòng \\
enwiki: event / edit / trang mới / ẩn danh / bot & 5 / 5 / 0 / 1 / 1 \\
viwiki: event / edit / trang mới / ẩn danh / bot & 1 / 0 / 1 / 0 / 0 \\
\bottomrule
\end{tabularx}
\end{table}

\begin{table}[H]
\centering\small
\caption{Các lần kiểm chứng}\label{tab:kiemchung}
\begin{tabularx}{\linewidth}{@{}llX@{}}
\toprule
\textbf{Ngày} & \textbf{Máy} & \textbf{Kết quả} \\ \midrule
11/10/2026 & Apple M3, ARM64 & 84 test nhẹ đạt; 16 test Spark đạt; kiểm thử end-to-end đạt (Trọng) \\
15/10/2026 & Windows 11, x86\_64 & CP0 chạy lại, số khớp Bảng~\ref{tab:cp0}; phát lại lần hai không đổi số; 33 test API đạt, gồm 4 test trên Cassandra thật (Nhân) \\
\bottomrule
\end{tabularx}
\end{table}
\dukiensua{Mỗi thành viên chạy lại CP0 trên máy mình thì thêm một dòng (CP1).}

\begin{ghichu}{Bài học khi chạy trên Windows}
Hai lỗi chỉ xuất hiện trên Windows. (1) Git trên Windows mặc định đổi ký tự xuống dòng của file sang kiểu Windows (CRLF);
script tạo topic chạy trong container Linux vì thế đọc sai và dừng. Cách sửa: khai báo cho Git luôn giữ kiểu Linux (LF)
với mọi script shell. (2) Thư viện Cassandra cho Python cần một module đã bị gỡ khỏi Python 3.12; container dùng
Python 3.11 nên không bị, còn máy phát triển dùng 3.12 phải cài thêm gói thay thế. Bài học: mỗi thành viên nên chạy lại
toàn bộ hệ thống trên máy mình sớm (mốc CP1), vì lỗi môi trường không lộ ra khi chỉ chạy trên một máy.
\end{ghichu}

\section{Xu hướng và ưu tiên xem xét trên fixture vàng}\chualam
\section{Độ trễ end-to-end}\chualam
\section{Khởi động lại, retry và dữ liệu đến muộn}\chualam
\section{Chạy LIVE liên tục}\chualam
\section{Tải và dung lượng}\chualam

% ============================================================================
% CHƯƠNG 6: KẾT LUẬN
% ============================================================================
\chapter{Kết luận và hướng phát triển}\label{chap:ketluan}

\section{Kết quả đạt được}
Đến ngày 15/10/2026, khung end-to-end đã chạy thật ở chế độ REPLAY: Kafka nhận sự kiện, Spark kiểm tra, loại trùng,
tổng hợp theo event time và ghi ba tầng Parquet, Cassandra phục vụ số tạm tính qua API, dashboard hiển thị trang
Overview. Kết quả khớp bộ kỳ vọng độc lập trên hai máy khác kiến trúc. Các nguyên tắc đúng khi retry (upsert tổng tuyệt
đối), riêng tư (không có định danh ngoài Bronze) và không dùng dữ liệu giả đã có test bảo vệ.
\dukiensua{Viết lại khi chốt CP4.}

\section{Hạn chế}
\begin{itemize}
  \item Một node, chưa có xác thực, phân quyền hay độ sẵn sàng cao.
  \item Không có giao dịch chung giữa các sink; tính đúng dựa vào upsert theo khóa ổn định.
  \item Bronze giữ payload gốc có thể chứa tên người dùng, IP; cần chính sách lưu giữ.
\end{itemize}

\section{Hướng phát triển}
Các mục dưới đây đã có trong kế hoạch nhưng chưa thực hiện; mỗi mục được viết đầy đủ khi hoàn thành.

\subsection{Collector LIVE với nối lại và đếm lỗi (W2-THANH-01)}\chualam
\subsection{Cửa sổ trượt, đặc trưng xu hướng và điểm rủi ro (W2-LINH-01)}\chualam
\subsection{API xu hướng, chỉnh sửa đáng ngờ và chi tiết sự kiện (W2-NHAN-01)}\chualam
\subsection{Hoàn thiện bốn trang dashboard (W2-MANH-01)}\chualam
\subsection{Profile HDFS và tích hợp ba tầng (W2-TRONG-01, W3-LINH-01)}\chualam
\subsection{Retry, latest window, TTL và kiểm thử dung lượng (W3-NHAN-01)}\chualam
\subsection{Khung Kubernetes và monitoring (W3-TRONG-01)}\chualam
\subsection{Đóng băng, clone sạch và demo (tuần 4)}\chualam

\appendix
\chapter{Bảng tra mã nguồn}\label{app:ma}
Phụ lục này nối các khái niệm trong báo cáo với vị trí mã nguồn, để người đọc tìm code mà không phải đọc tên file trong
phần nội dung. Đường dẫn tính từ thư mục gốc của repo.

\begin{table}[H]
\centering\small
\caption{Khái niệm và vị trí mã nguồn}\label{tab:trama}
\begin{tabularx}{\linewidth}{@{}Xl@{}}
\toprule
\textbf{Khái niệm} & \textbf{Vị trí} \\ \midrule
Cấu hình, kiểm tra tham số khởi động & \code{shared/wikipulse\_core/config.py} \\
Lớp truy cập Cassandra, kiểm tra dòng trước khi ghi & \code{shared/wikipulse\_core/serving.py} \\
Tạo topic Kafka & \code{scripts/create\_topics.sh} \\
Phát lại dữ liệu mẫu & \code{replay/replay\_producer.py} \\
Dữ liệu mẫu và kết quả kỳ vọng & \code{tests/fixtures/} \\
Tiến trình Spark và năm truy vấn & \code{processing/spark/jobs/streaming\_job.py} \\
Kiểm tra, chuẩn hóa, loại trùng, đếm & \code{processing/spark/wikipulse/transformations.py} \\
Ghi kết quả sang Cassandra & \code{processing/spark/wikipulse/sinks.py} \\
Thư mục theo mã lần chạy & \code{processing/spark/wikipulse/paths.py} \\
Lược đồ Cassandra & \code{storage/cassandra/schema.cql} \\
API & \code{api/app/main.py}, \code{api/app/models.py} \\
Dashboard & \code{dashboard/app.py}, \code{dashboard/views/} \\
Contract giữa các module & \code{docs/contracts.md} \\
Bảng việc và trạng thái & \code{docs/checkpoints.md} \\
\bottomrule
\end{tabularx}
\end{table}

\begin{table}[H]
\centering\small
\caption{Lệnh thường dùng (chạy ở thư mục gốc repo)}\label{tab:lenh}
\begin{tabularx}{\linewidth}{@{}lX@{}}
\toprule
\textbf{Lệnh} & \textbf{Tác dụng} \\ \midrule
\code{make install} & Tạo môi trường Python và cài thư viện \\
\code{make up} & Dựng và khởi động toàn bộ hệ thống \\
\code{make replay} & Phát lại dữ liệu mẫu vào Kafka \\
\code{make down} & Dừng hệ thống, giữ dữ liệu \\
\code{make test} & Chạy kiểm thử nhẹ \\
\code{make test-spark} & Chạy kiểm thử Spark trong container \\
\code{pytest -m integration api/tests} & Kiểm thử trên Cassandra thật (cần đặt địa chỉ Cassandra qua biến môi trường) \\
\code{make smoke} & Kiểm thử end-to-end trên hệ thống riêng, tự dọn sau khi chạy \\
\bottomrule
\end{tabularx}
\end{table}

\clearpage\phantomsection
\addcontentsline{toc}{chapter}{Tài liệu tham khảo}
\bibliographystyle{plain}
\bibliography{references}

\end{document}
